\section*{Capítulo \ref{ch:b2:metric} --- Topología de los espacios métricos}

\begin{solution}{exo:b2:metric:1}
$\delta$: la simetría y la separación son claras; desigualdad
triangular: $\min(1, u + v) \leq \min(1,u) + \min(1,v)$ para $u, v
\geq 0$ (si alguno de los dos mínimos vale $1$, el miembro derecho es
$\geq 1$; en caso contrario vale $u + v$). $d$: es el retroceso de
$\abs{\cdot}$ por la aplicación inyectiva $\arctan$, y los tres
axiomas se transfieren.

Convergencia: para $\delta$, $\delta(x_n, x) \to 0 \iff \abs{x_n - x}
\to 0$ (para valores pequeños ambas distancias coinciden): las mismas
sucesiones convergentes de siempre. Para $d$: $d(x_n, x) \to 0 \iff
\arctan x_n \to \arctan x \iff x_n \to x$ (\omterm{def:b2:metric:continuity}{continuidad} y monotonía
estricta de $\arctan$ y de su inversa en los rangos pertinentes): de
nuevo, la convergencia usual.

$(\R, d)$ \emph{no} es \omterm{def:b2:metric:complete}{completo}: $x_n = n$ cumple $d(x_p, x_q) =
\abs{\arctan p - \arctan q} \to 0$ (ambos tienden a $\frac\pi2$),
luego es de Cauchy; pero $(x_n)$ no converge para $d$ (su $d$-límite
sería un límite ordinario). La \omterm{def:b2:metric:complete}{completitud} es una propiedad de la
\emph{distancia}, no solo de las sucesiones convergentes.
\end{solution}

\begin{solution}{exo:b2:metric:2}
Convergente $\Rightarrow$ de Cauchy: $d(x_p, x_q) \leq d(x_p, \ell) +
d(\ell, x_q)$. Acotada: a partir de $N$, $d(x_n, \ell) \leq 1$; y los
finitos primeros términos también caben en algún radio.

De Cauchy con subsucesión convergente $x_{\varphi(n)} \to \ell$: dado
$\varepsilon$, para $n$ grande, $d(x_n, \ell) \leq d(x_n,
x_{\varphi(n)}) + d(x_{\varphi(n)}, \ell) \leq 2\varepsilon$ (el
primer término, por ser de Cauchy, ya que $\varphi(n) \geq n$).

\omterm{def:b2:metric:compact}{Compacto} $\Rightarrow$ \omterm{def:b2:metric:complete}{completo}: una sucesión de Cauchy tiene una
subsucesión convergente (\omterm{def:b2:metric:compact}{compacidad}) y por tanto converge.
\end{solution}

\begin{solution}{exo:b2:metric:3}
$d_\infty(f, g) = \sup_{\intcc{0}{1}} \abs{x - x^2} = \frac14$
(máximo de $x - x^2$ en $x = \frac12$).

$\overline B(0, 1) = \{f : \sup\abs f \leq 1\}$: las funciones
\omterm{def:b2:metric:continuity}{continuas} con valores en $\intcc{-1}{1}$.

$\{f : f(0) = 0\}$ es la imagen recíproca de $\{0\}$ por la
\emph{evaluación} $f \mapsto f(0)$, que es $1$-lipschitziana
($\abs{f(0) - g(0)} \leq d_\infty(f,g)$) y por tanto \omterm{def:b2:metric:continuity}{continua}: el
conjunto es cerrado (\cref{thm:b2:metric:globalcontinuity}).
Análogamente, $\{f : f(0) > 0\}$ es la imagen recíproca del \omterm{def:b2:metric:topology}{abierto}
$\intoo{0}{+\infty}$, luego es \omterm{def:b2:metric:topology}{abierto}.
\end{solution}

\begin{solution}{exo:b2:metric:4}
\emph{\omterm{def:b2:metric:complete}{Completo}:} una sucesión de Cauchy con $\varepsilon = \frac12$
es constante a partir de un índice y, por tanto, convergente.

\emph{\omterm{def:b2:metric:compact}{Compacto} si y solo si finito:} si $X$ es finito, toda sucesión
toma algún valor infinitas veces (subsucesión constante). Si $X$ es
infinito, una sucesión de puntos distintos dos a dos tiene todas sus
distancias mutuas iguales a $1$: no hay subsucesión de Cauchy y, por
tanto, tampoco convergente.

\emph{Subconjuntos \omterm{def:b2:metric:connected}{conexos}:} los conjuntos unitarios (y
$\emptyset$). Todo $A$ con dos puntos $x \neq y$ se parte como $\{x\}
\cup (A \setminus\{x\})$, ambos \omterm{def:b2:metric:topology}{abiertos} en $A$ (todo subconjunto de
un espacio discreto es \omterm{def:b2:metric:topology}{abierto}, pues las bolas de radio $\frac12$ son
unitarias).
\end{solution}

\begin{solution}{exo:b2:metric:5}
La función $g(x) = d(x, f(x))$ es \omterm{def:b2:metric:continuity}{continua} sobre el \omterm{def:b2:metric:compact}{compacto} $X$
(pues $\abs{g(x) - g(y)} \leq 2d(x,y)$, por dos desigualdades
triangulares), así que alcanza su mínimo en cierto $a$
(\cref{thm:b2:metric:compactprops}). Si $f(a) \neq a$:
\[
g\bigl(f(a)\bigr) = d\bigl(f(a), f(f(a))\bigr) < d(a, f(a)) = g(a),
\]
en contra de la minimalidad. Luego $f(a) = a$; la unicidad es la de
siempre (dos puntos fijos $a \neq b$ darían $d(a,b) = d(f(a), f(b)) <
d(a,b)$).

Ejemplo no \omterm{def:b2:metric:compact}{compacto}: $f(x) = x + \frac1x$ sobre
$\intco{1}{+\infty}$: $\abs{f(x) - f(y)} = \abs{x - y}\,\abs{1 -
\frac{1}{xy}} < \abs{x - y}$ para $x \neq y$ (ya que $xy > 1$) y, sin
embargo, $f(x) > x$ en todas partes.
\end{solution}

\begin{solution}{exo:b2:metric:6}
Tomemos $x_n \in K_n$. Todos los términos a partir del índice $n$
están en $K_n$; en particular, toda la sucesión está en el
\omterm{def:b2:metric:compact}{compacto} $K_0$: extraigamos $x_{\varphi(k)} \to \ell$. Para cada $n$
fijo, los términos $x_{\varphi(k)}$ con $\varphi(k) \geq n$ están en
el \emph{cerrado} $K_n$, luego el límite $\ell \in K_n$. Por tanto
$\ell \in \bigcap K_n$.

Contraejemplo con cerrados: $F_n = \intco{n}{+\infty}$ en $\R$:
encajados, cerrados, no vacíos y de intersección vacía.
\end{solution}

\begin{solution}{exo:b2:metric:7}
Que $f^{-1}$ sea \omterm{def:b2:metric:continuity}{continua} significa que las imágenes $f(F)$ de los
cerrados $F \subseteq K$ son cerradas (las imágenes recíprocas por
$f^{-1}$ son las imágenes por $f$). Un cerrado $F$ del \omterm{def:b2:metric:compact}{compacto} $K$
es \omterm{def:b2:metric:compact}{compacto} (\cref{thm:b2:metric:compactprops}~(1)); su imagen
\omterm{def:b2:metric:continuity}{continua} $f(F)$ es \omterm{def:b2:metric:compact}{compacta} y, por tanto, cerrada. Luego $f^{-1}$
es \omterm{def:b2:metric:continuity}{continua} y $f$ es un homeomorfismo.

Contraejemplo: $\gamma(t) = (\cos t, \sin t)$ de $\intco{0}{2\pi}$
(no \omterm{def:b2:metric:compact}{compacto}) sobre la circunferencia unidad es una biyección
\omterm{def:b2:metric:continuity}{continua}, pero $\gamma^{-1}$ es discontinua en $(1,0)$: los puntos
de la circunferencia justo por debajo del eje tienen parámetros
cercanos a $2\pi$, no a $0$.
\end{solution}

\begin{solution}{exo:b2:metric:8}
$C = \bigcap_n C_n$, donde cada $C_n$ (unión de $2^n$ intervalos
cerrados de longitud $3^{-n}$) es cerrado: $C$ es cerrado y acotado
en $\R$ y, por tanto, \omterm{def:b2:metric:compact}{compacto}
(\cref{thm:b2:metric:compactprops}~(2)).

Interior vacío: $C$ no contiene ningún intervalo de longitud $>
3^{-n}$ (está dentro de $C_n$, cuyas componentes tienen esa
longitud), y esto para todo $n$.

Cardinal: los puntos de $C$ son exactamente los reales
$\sum_{n\geq1} a_n 3^{-n}$ con cifras $a_n \in \{0, 2\}$ (en cada
etapa, el tercio central suprimido elimina la cifra $1$); la
aplicación $(a_n) \mapsto \sum a_n 3^{-n}$ es una biyección de
$\{0,2\}^{\N^*}$ sobre $C$ (la inyectividad, como en el~\cref{exo:b2:structures:3}). Así pues, $C$ es \omterm{def:b2:structures:countable}{equipotente} a
$\{0,1\}^{\N}$: no \omterm{def:b2:structures:countable}{numerable}, pese a tener ``longitud nula''.
\end{solution}

\begin{solution}{exo:b2:metric:9}
Supongamos $a \notin f(X)$. La imagen $f(X)$ es \omterm{def:b2:metric:compact}{compacta} (imagen
\omterm{def:b2:metric:continuity}{continua}) y, por tanto, cerrada; luego
\[
\varepsilon = d\bigl(a, f(X)\bigr)
= \inf_{y \in f(X)} d(a, y) > 0
\]
(el ínfimo de una función \omterm{def:b2:metric:continuity}{continua} sobre un \omterm{def:b2:metric:compact}{compacto} se alcanza; si
fuera $0$, $a$ sería adherente al cerrado $f(X)$ y por tanto
pertenecería a él).

Consideremos la órbita $x_n = f^n(a)$ (con $x_0 = a$). Para $p < q$:
\[
d(x_p, x_q) = d\bigl(f^p(a), f^p(f^{q-p}(a))\bigr)
= d\bigl(a, f^{q-p}(a)\bigr) \geq \varepsilon,
\]
(la isometría iterada $p$ veces; y $f^{q-p}(a) \in f(X)$ porque $q -
p \geq 1$). Una sucesión con distancias mutuas $\geq \varepsilon$ no
tiene subsucesión convergente, en contra de la \omterm{def:b2:metric:compact}{compacidad}. Por
tanto $f(X) = X$.
\end{solution}

\begin{solution}{exo:b2:metric:10}
Sean $A, B \in GL_n(\C)$ y $p(z) = \det\bigl((1-z)A + zB\bigr)$: un
polinomio en $z$ (cada entrada es afín en $z$ y el \omterm{def:b2:linalg:det}{determinante} es
polinómico en las entradas). Como $p(0) = \det A \neq 0$, $p$ no es
idénticamente nulo y tiene un número finito de raíces $z_1, \dots,
z_m$ (ninguna igual a $0$ ni a $1$, pues $p(1) = \det B \neq 0$). El
plano $\C$ menos un número finito de puntos es \omterm{def:b2:metric:connected}{conexo} por caminos:
existe un camino de $0$ a $1$ que evita los $z_i$ (tómese una línea
quebrada que pase por un punto alejado de todas las raíces, o un arco
de circunferencia; los obstáculos son finitos). A lo largo de un
camino así, $\gamma$, la aplicación $t \mapsto (1 - \gamma(t))A +
\gamma(t)B$ es un camino \omterm{def:b2:metric:continuity}{continuo} \emph{dentro} de $GL_n(\C)$ que va
de $A$ a $B$ (el \omterm{def:b2:linalg:det}{determinante} no se anula en él). Así pues,
$GL_n(\C)$ es \omterm{def:b2:metric:connected}{conexo} por caminos, al contrario que su primo real
(\cref{ex:b2:metric:glnr}): el plano complejo tiene sitio para
rodear los obstáculos.
\end{solution}

\begin{solution}{exo:b2:metric:11}
\emph{Lipschitz:} para $a \in A$, $d(x, a) \leq d(x, y) + d(y,
a)$; tomando el ínfimo sobre $a$: $d(x, A) \leq d(x, y) + d(y, A)$,
e intercambiando $x$ e $y$: $\abs{d(x,A) - d(y,A)} \leq d(x,y)$.

\emph{Anulación:} $d(x, A) = 0$ si y solo si existen $a_n \in A$ con
$d(x, a_n) \to 0$, si y solo si $x$ es límite de puntos de $A$, si y
solo si $x \in \overline A$.

\emph{El interruptor:} para cerrados disjuntos $A, B$, el
denominador $d(x,A) + d(x,B)$ no se anula nunca (ello obligaría a $x
\in \overline A \cap \overline B = A \cap B = \emptyset$), de modo
que $\varphi$ está bien definida y es \omterm{def:b2:metric:continuity}{continua} como cociente de
funciones \omterm{def:b2:metric:continuity}{continuas} con denominador no nulo. $\varphi(x) = 0$ si y
solo si $d(x, A) = 0$, si y solo si $x \in A$; $\varphi(x) = 1$ si y
solo si $d(x, B) = 0$, si y solo si $x \in B$; y $0 \leq \varphi \leq
1$ en todas partes.
\end{solution}

\begin{solution}{exo:b2:metric:12}
Sea $B(x_0, r_0)$ una bola cualquiera; hallemos en ella un punto de
$\bigcap U_n$. Como $U_1$ es denso y \omterm{def:b2:metric:topology}{abierto}, $U_1 \cap B(x_0, r_0)$
es no vacío y \omterm{def:b2:metric:topology}{abierto}: contiene una bola cerrada $\overline
B(x_1, r_1)$ con $0 < r_1 \leq \frac{r_0}2$ (redúzcase el radio).
Por inducción, $U_{n+1} \cap B(x_n, r_n)$ es \omterm{def:b2:metric:topology}{abierto} no vacío:
tómese $\overline B(x_{n+1}, r_{n+1}) \subseteq U_{n+1} \cap
B(x_n, r_n)$ con $r_{n+1} \leq \frac{r_n}2$. Para $p, q \geq n$,
tanto $x_p$ como $x_q$ están en $B(x_n, r_n)$ con $r_n \leq
2^{-n}r_0$: la sucesión es de Cauchy y converge a cierto $\ell$
por \omterm{def:b2:metric:complete}{completitud}. Para cada $n$, la cola de la sucesión está en la
bola \emph{cerrada} $\overline B(x_{n+1}, r_{n+1}) \subseteq
U_{n+1} \cap B(x_0, r_0)$, luego $\ell \in U_{n+1}$ para todo $n$
y $\ell \in \overline B(x_1, r_1) \subseteq B(x_0, r_0)$. Por
tanto $\bigcap_n U_n$ corta a toda bola: es denso.

\emph{Aplicación:} si $\R = \bigcup_n F_n$ con $F_n$ cerrados de
interior vacío, entonces los $U_n = \R \setminus F_n$ son densos
($\overline{U_n} = \R$ si y solo si $F_n$ tiene interior vacío) y
\omterm{def:b2:metric:topology}{abiertos}, y Baire da un punto de $\bigcap U_n = \R \setminus
\bigcup F_n$: contradicción. En particular, $\R \neq \bigcup_{x
\in D} \{x\}$ para un $D$ \omterm{def:b2:structures:countable}{numerable} (los conjuntos unitarios son
cerrados de interior vacío): $\R$ no es \omterm{def:b2:structures:countable}{numerable}, el~\cref{thm:b2:structures:cantor} por otro camino.
\end{solution}

\begin{solution}{pb:b2:metric:1}
\textbf{1.} Sea $F$ cerrado en el \omterm{def:b2:metric:complete}{completo} $X$ y $(x_n) \subseteq
F$ de Cauchy: converge en $X$ a cierto $\ell$, y $\ell \in F$
porque $F$ es cerrado (los límites de sucesiones de $F$ se quedan
en $\overline F = F$): $F$ es \omterm{def:b2:metric:complete}{completo}. Recíprocamente, sea $A
\subseteq X$ \omterm{def:b2:metric:complete}{completo} y $x \in \overline A$: alguna sucesión de
$A$ converge a $x$; es de Cauchy, luego converge \emph{en} $A$; y
como los límites son únicos, $x \in A$, es decir, $A$ es cerrado.
Puesto que $\bigl(C(I), d_\infty\bigr)$ es \omterm{def:b2:metric:complete}{completo}
(\cref{thm:b2:metric:rncomplete}), sus cerrados son \omterm{def:b2:metric:complete}{completos}.

\textbf{2.} Para $y, z \in C(I)$ y $t \in I$:
\[
\Bigl|\int_{t_0}^{t}y - \int_{t_0}^{t}z\Bigr| \leq
\abs{t - t_0}\,\sup_I\abs{y - z} \leq h\,d_\infty(y, z),
\]
y se toma el supremo sobre $t$.

\textbf{3.} Usando las dos ecuaciones de punto fijo y la
desigualdad triangular:
\[
d(\ell_g, \ell_{\widetilde g})
= d\bigl(g(\ell_g), \widetilde g(\ell_{\widetilde g})\bigr)
\leq d\bigl(g(\ell_g), g(\ell_{\widetilde g})\bigr) +
d\bigl(g(\ell_{\widetilde g}), \widetilde g(\ell_{\widetilde
g})\bigr)
\leq k\,d(\ell_g, \ell_{\widetilde g}) +
d\bigl(g(\ell_{\widetilde g}), \widetilde g(\ell_{\widetilde
g})\bigr),
\]
y se despeja $d(\ell_g, \ell_{\widetilde g})$ (el coeficiente $1 -
k$ es positivo). La segunda desigualdad acota la separación
evaluada por la uniforme.

\textbf{4.} $g^m$ es una contracción de un espacio \omterm{def:b2:metric:complete}{completo} no
vacío: tiene un único punto fijo $\ell$
(\cref{thm:b2:metric:banach}). Entonces $g^m(g(\ell)) =
g(g^m(\ell)) = g(\ell)$, así que $g(\ell)$ es punto fijo de $g^m$
y, por unicidad, $g(\ell) = \ell$. Todo punto fijo de $g$ lo es de
$g^m$: de ahí la unicidad para $g$. Órbitas: fijemos $r \in \{0,
\dots, m-1\}$; la subsucesión $(x_{qm + r})_q$ es la órbita de
$g^m$ que arranca en $x_r$, luego converge a $\ell$ cuando $q \to
\infty$ (Banach de nuevo). Las $m$ subsucesiones convergen al
mismo $\ell$, de donde $x_n \to \ell$: dado $\varepsilon$, cada
clase de restos acaba a distancia $< \varepsilon$, y las clases
son finitas en número.

\textbf{5.} Si $y$ es \omterm{def:b2:metric:continuity}{continua} con valores en $\intcc{y_0 -
b}{y_0 + b}$, el integrando $s \mapsto f(s, y(s))$ es \omterm{def:b2:metric:continuity}{continuo}
sobre $I$ (composición), de modo que el miembro derecho es de
clase $C^1$ con derivada $f(t, y(t))$ (teorema fundamental del
cálculo, volumen del primer año). Si $y$ satisface la ecuación
integral, es esa función $C^1$, cumple $y(t_0) = y_0$ e $y' = f(t,
y)$. Recíprocamente, integrando $y' = f(s, y(s))$ de $t_0$ a $t$
se obtiene la ecuación integral.

\textbf{6.} $X_h$ contiene la función constante $y_0$; es cerrado
por ser la imagen recíproca de $\intcc{0}{b}$ por la aplicación
\omterm{def:b2:metric:continuity}{continua} $y \mapsto d_\infty(y, y_0)$ (las distancias son
$1$-lipschitzianas) y, por tanto, \omterm{def:b2:metric:complete}{completo} por la pregunta 1.
Estabilidad: para $y \in X_h$ y $t \in I$,
\[
\abs{T(y)(t) - y_0} = \Bigl|\int_{t_0}^{t}f(s, y(s))\dd s\Bigr|
\leq M\abs{t - t_0} \leq Mh \leq b ,
\]
donde el último paso usa $h \leq b/M$ (o $M = 0$, caso trivial). Y
$T(y)$ es \omterm{def:b2:metric:continuity}{continua} (incluso $C^1$, pregunta 5): $T(y) \in X_h$.

\textbf{7.} Para $t \in I$:
\[
\abs{T(y)(t) - T(z)(t)} \leq \int_{t_0}^{t}\abs{f(s, y(s)) -
f(s, z(s))}\,\abs{\dd s} \leq L\abs{t - t_0}\,d_\infty(y,z)
\leq Lh\,d_\infty(y,z).
\]
Si $Lh < 1$, $T$ es una contracción del \omterm{def:b2:metric:complete}{completo} no vacío $X_h$ y
Banach da un único punto fijo que, por la pregunta 5, es la única
solución.

\textbf{8.} Por inducción; el caso $n = 1$ es la desigualdad
central de la pregunta 7. Supuesta la cota para $n$, y para $t
\geq t_0$ (el caso $t \leq t_0$ es simétrico):
\[
\abs{T^{n+1}(y)(t) - T^{n+1}(z)(t)}
\leq L\int_{t_0}^{t}\abs{T^n(y)(s) - T^n(z)(s)}\dd s
\leq L\int_{t_0}^{t}\frac{L^n(s - t_0)^n}{n!}\dd s\;
d_\infty(y,z),
\]
y la integral vale $\frac{L^n(t - t_0)^{n+1}}{(n+1)!}$: la cota
para $n + 1$. Por tanto $d_\infty(T^n y, T^n z) \leq
\frac{(Lh)^n}{n!}d_\infty(y, z)$, y $\frac{(Lh)^n}{n!} \to 0$
(la serie exponencial converge): alguna $T^m$ es una contracción.
La pregunta 4 se aplica sobre el \omterm{def:b2:metric:complete}{completo} $X_h$: $T$ tiene un
único punto fijo, es decir, el problema de Cauchy tiene
exactamente una solución sobre $I$ con valores en $\intcc{y_0 -
b}{y_0 + b}$.

\textbf{9.} Sea $y$ solución del problema sobre $I$ y supongamos
que el conjunto $E = \{t \in I, t > t_0 : \abs{y(t) - y_0} > b\}$
es no vacío (el lado $t < t_0$ es simétrico); sea $\tau = \inf E$.
Por \omterm{def:b2:metric:continuity}{continuidad}, $\abs{y(s) - y_0} \leq b$ para $s \in
\intcc{t_0}{\tau}$, de modo que allí la gráfica está en $R$, la
ecuación integral vale hasta $\tau$ y
\[
\abs{y(\tau) - y_0} = \Bigl|\int_{t_0}^{\tau}f\bigl(s,
y(s)\bigr)\dd s\Bigr| \leq M(\tau - t_0) \leq Mh \leq b .
\]
Si $\tau < t_0 + h$, los puntos de $E$ arbitrariamente próximos a
$\tau$ por la derecha dan, por \omterm{def:b2:metric:continuity}{continuidad}, $\abs{y(\tau) - y_0}
\geq b$ y por tanto $= b$; pero entonces la fórmula anterior
obliga a $M(\tau - t_0) = Mh$, es decir, $\tau = t_0 + h$:
contradicción. Así pues, $\tau = t_0 + h$, $E \subseteq \{t_0 +
h\}$, y la fórmula (en $\tau = t_0 + h$) da $\abs{y(t_0 + h) -
y_0} \leq b$, en contra de la pertenencia a $E$. Luego $E =
\emptyset$: toda solución sobre $I$ permanece en la banda, es
punto fijo de $T$ en $X_h$, y la unicidad es incondicional.

\textbf{10.} $T(y)(t) = 1 + \int_0^t y$. Partiendo de $y^{(0)}
\equiv 1$:
\[
y^{(1)}(t) = 1 + t,\quad
y^{(2)}(t) = 1 + t + \frac{t^2}2,\quad\dots\quad
y^{(n)}(t) = \sum_{k=0}^{n}\frac{t^k}{k!}
\]
(por inducción: integrar la suma parcial añade el término
siguiente). Son las sumas parciales de Taylor de $\eu^{\,t}$;
sobre cualquier $I$ acotado convergen uniformemente a $\eu^{\,t}$
(la cola está dominada por la serie numérica convergente $\sum
h^k/k!$), que es en efecto la única solución.

\textbf{11.} $y \equiv 0$ es solución. Para $y_c$: es de clase
$C^1$ (ambos trozos lo son y en $t = c$ las derivadas coinciden:
$0$ y $2(t - c) \to 0$), y para $t > c$ se tiene $y_c' = 2(t - c)
= 2\sqrt{(t-c)^2} = 2\sqrt{\abs{y_c}}$; para $t \leq c$ ambos
miembros se anulan. Así pues, el problema de Cauchy $y(0) = 0$
tiene infinitas soluciones ($c \geq 0$ arbitrario, más $y \equiv
0$). El campo $\varphi(y) = 2\sqrt{\abs y}$ no es \omterm{def:b2:metric:continuity}{lipschitziano}
cerca de $0$: $\frac{\varphi(y) - \varphi(0)}{y - 0} =
\frac{2}{\sqrt y} \to +\infty$ cuando $y \to 0^+$, de modo que
ninguna constante $L$ sirve en ningún entorno de $0$, justo donde
todas las soluciones se ramifican.

\textbf{12.} Separando variables (o comprobando directamente), la
única solución local es $y(t) = \frac1{1 - t}$, definida sobre
$\intoo{-\infty}{1}$ y con explosión en $t = 1$. Para el
rectángulo $\intcc{-a}{a} \times \intcc{1 - b}{1 + b}$: $M = \sup
y^2 = (1 + b)^2$, luego la semianchura certificada es $h =
\min\bigl(a, \frac{b}{(1+b)^2}\bigr)$. Maximizando
$\frac{b}{(1+b)^2}$: la derivada se anula en $b = 1$, con valor
$\frac14$. Así pues, el teorema solo garantiza vida sobre
$\intcc{-\frac14}{\frac14}$, correctamente menos que la duración
verdadera $1$ hacia delante e infinitamente menos hacia atrás: el
teorema es local por naturaleza, y la explosión muestra que no
puede ser de otro modo.

\textbf{13.} $g$ envía $\intcc12$ en sí mismo: $g$ decrece sobre
$\intcc{1}{\sqrt2}$ y crece después (estúdiese $g'(x) = \frac12 -
\frac1{x^2}$), con $g(1) = g(2) = \frac32$ y mínimo $g(\sqrt2) =
\sqrt2 > 1$: $g(\intcc12) \subseteq \intcc{\sqrt2}{\frac32}
\subseteq \intcc12$; y $g$ envía racionales a racionales.
Contracción: $\abs{g'(x)} = \abs{\frac12 - \frac1{x^2}} \leq
\frac12$ sobre $\intcc12$ (pues $\frac1{x^2} \in
\intcc{\frac14}{1}$), de modo que la desigualdad del valor medio
da $\abs{g(x) - g(y)} \leq \frac12\abs{x - y}$. Un punto fijo
cumpliría $\frac x2 = \frac1x$, es decir, $x^2 = 2$: imposible en
$\Q$. La hipótesis que falla es la \omterm{def:b2:metric:complete}{completitud} de $X$ ($\Q \cap
\intcc12$ no es \omterm{def:b2:metric:complete}{completo}); en el completado $\intcc12$ el punto
fijo es $\sqrt2$: el teorema de Banach ejecutado sobre los
racionales \emph{crea} el irracional.

\textbf{14.} A posteriori: $d(x_n, \ell) \leq d(x_n, x_{n+1}) +
d(x_{n+1}, \ell) \leq k\,d(x_{n-1}, x_n) + k\,d(x_n, \ell)$, de
donde $d(x_n, \ell) \leq \frac{k}{1-k}d(x_n, x_{n-1})$. Herón
desde $x_0 = \frac32$: $x_1 = \frac{17}{12}$, $d(x_1, x_0) =
\frac1{12}$, $k = \frac12$; la cota a priori
$\frac{k^n}{1-k}d(x_1,x_0) = \frac{2^{-n+1}}{12}$ baja de
$10^{-6}$ por primera vez en $n = 18$. En realidad, $x_1 =
\frac{17}{12} \approx 1.41667$ (error $2.5\cdot10^{-3}$), $x_2 =
\frac{577}{408} \approx 1.4142157$ (error $2.1\cdot10^{-6}$) y
$x_3 \approx 1.41421356237469$ (error $1.6\cdot10^{-12}$): bastan
tres pasos. Cada paso de Herón \emph{eleva al cuadrado}
aproximadamente el error (convergencia cuadrática, fenómeno de
Newton: \cref{ch:b2:realfun}); la estimación de la contracción,
que solo lo reduce a la mitad, es honesta para el peor caso pero
aquí resulta pesimista.

\textbf{15.} Apliquemos la pregunta 3 con $g = T_{y_0}$ (una
contracción de razón $Lh < 1$) y $\widetilde g = T_{z_0}$, cuyo
punto fijo es $y[z_0]$. Para cualquier $y$,
\[
\abs{T_{y_0}(y)(t) - T_{z_0}(y)(t)} = \abs{y_0 - z_0},
\]
(las integrales son idénticas), de modo que $\sup_y
d_\infty(T_{y_0}(y), T_{z_0}(y)) = \abs{y_0 - z_0}$, y la
pregunta 3 da
\[
d_\infty(y[y_0], y[z_0]) \leq \frac{\abs{y_0 - z_0}}{1 - Lh} .
\]

\textbf{16.} En cada trozo, la pregunta 15 aplicada tomando como
datos iniciales los valores en el extremo izquierdo acota la
desviación en el extremo derecho:
\[
d_\infty \leq \frac{1}{1 - 1/2}\,(\text{desviación izquierda})
= 2\,(\text{desviación izquierda}) .
\] Por
inducción sobre los $m$ trozos, la desviación final es a lo sumo
$2^m\abs{y_0 - z_0}$, y la desviación uniforme sobre todo el
segmento obedece la misma cota (el supremo de cada trozo queda
controlado en su etapa). Con trozos de longitud $h \asymp
\frac1{2L}$, el factor es $2^m = 2^{\,\text{longitud}\cdot 2L}$:
exponencial en la longitud del intervalo, exactamente como predice
la cota $\eu^{L\abs{t-t_0}}$ de Gronwall, con mejores constantes.

\textbf{17.} Mismo esquema: para $y \in X_h$,
\[
\abs{T_f(y)(t) - T_g(y)(t)} \leq \int_{t_0}^t \abs{f(s,y(s)) -
g(s,y(s))}\,\abs{\dd s} \leq \varepsilon h ,
\]
de modo que la pregunta 3 (con $T_f$ como contracción y $T_g$ como
aplicación perturbada) da $d_\infty \leq \frac{\varepsilon h}{1 -
Lh}$.

\textbf{18.} Por la pregunta 17 aplicada a $f = f_\lambda$, $g =
f_\mu$: $d_\infty(y_\lambda, y_\mu) \leq \frac{Ch}{1 -
Lh}\abs{\lambda - \mu}$; la aplicación solución es
\omterm{def:b2:metric:continuity}{lipschitziana} de constante $\frac{Ch}{1-Lh}$.

\textbf{19.} Todos los argumentos usaron únicamente los axiomas
métricos, la \omterm{def:b2:metric:complete}{completitud} del escenario, la cota $\abs{\int} \leq
\int\abs{\cdot}$ y el carácter \omterm{def:b2:metric:continuity}{lipschitziano} de $f$, todo ello
disponible para funciones con valores en $\R^n$ y $d_\infty$
construida sobre cualquiera de las distancias equivalentes del~\cref{ex:b2:metric:examples} (\omterm{def:b2:metric:complete}{completo} por el~\cref{thm:b2:metric:rncomplete}). Para $y' = Ay$ con $A$
nilpotente: $y^{(0)} \equiv (c_1, c_2)$,
\[
y^{(1)}(t) = (c_1, c_2) + \int_0^t (c_2, 0)\,\dd s
= (c_1 + tc_2,\; c_2),
\]
y de nuevo $Ay^{(1)}(s) = (c_2, 0)$: luego $y^{(2)} = y^{(1)}$. La
iteración se estaciona a partir de $n = 1$ en la solución exacta
$y(t) = (c_1 + tc_2, c_2) = \eu^{tA}y(0)$: la nilpotencia trunca
la serie exponencial, y Picard se da cuenta.

\textbf{20.} Fijemos $y \in X$ y sea $g_y(x) = y - \eta(x)$: una
contracción de razón $k$ del \omterm{def:b2:metric:complete}{completo} $X$, luego hay exactamente
un $x$ con $x + \eta(x) = y$, es decir, la aplicación $\Phi =
\mathrm{id} + \eta$ es biyectiva. Inversa \omterm{def:b2:metric:continuity}{lipschitziana}: si
$\Phi(x) = y$ y $\Phi(x') = y'$, entonces
\[
d(x, x') \leq d(y, y') + d\bigl(\eta(x), \eta(x')\bigr)
\leq d(y, y') + k\,d(x, x'),
\]
luego $d(x, x') \leq \frac{1}{1-k}d(y, y')$. (Las distancias
proceden aquí de la estructura de norma, de modo que $d(a - c, b -
c) = d(a, b)$, hecho que usa la primera desigualdad.)

\textbf{21.} $g(x) = m + e\sin x$ es $e$-lipschitziana sobre el
\omterm{def:b2:metric:complete}{completo} $\R$ (desigualdad del valor medio: $\abs{g'} =
\abs{e\cos x} \leq e < 1$), así que Banach da una única solución y
la convergencia global de la iteración. Para $e = \frac12$, $m =
1$ y $x_0 = 1$: $x_1 = 1 + \frac12\sin 1 \approx 1.42074$, $d(x_1,
x_0) \approx 0.4207$, y la cota a priori
$\frac{(1/2)^n}{1/2}\cdot 0.4207 \leq 10^{-3}$ se cumple por
primera vez en $n = 10$: diez iteraciones certificadas (el valor
verdadero $x \approx 1.4987$ se alcanza de hecho con precisión
$10^{-3}$ hacia $n = 5$).

\textbf{22.} Usemos la descripción por cifras
(\cref{exo:b2:metric:8}): $C$ es el conjunto de las sumas
$\sum_{n\geq1}a_n3^{-n}$ con $a_n \in \{0, 2\}$. Entonces $S_1(C)
= \{x/3 : x \in C\}$ es el subconjunto con $a_1 = 0$, y $S_2(C) =
\{x/3 + 2/3\}$ el subconjunto con $a_1 = 2$: su unión, al recorrer
libremente $a_1$, es exactamente $C$. Unicidad en una frase: sobre
el espacio de los \omterm{def:b2:metric:compact}{compactos} no vacíos de $\intcc01$ dotado de la
distancia de Hausdorff, $A \mapsto S_1(A) \cup S_2(A)$ es una
contracción de razón $\frac13$ de un espacio \omterm{def:b2:metric:complete}{completo}, luego
Banach solo permite un conjunto fijo; el volumen del tercer año
formaliza la métrica de Hausdorff y este argumento.

\textbf{23.} Sea $Z = \{t \in J : y(t) = z(t)\}$: es no vacío
($t_0 \in Z$) y cerrado en $J$ (igualador de dos aplicaciones
\omterm{def:b2:metric:continuity}{continuas}: imagen recíproca de $\{0\}$ por $y - z$). \omterm{def:b2:metric:topology}{Abierto}:
si $t_1 \in Z$, apliquemos el teorema local (pregunta 8) en el
punto $(t_1, y(t_1))$, dentro de un rectángulo donde $f$ sea
\omterm{def:b2:metric:continuity}{lipschitziana}: sobre un intervalo pequeño en torno a $t_1$, tanto
$y$ como $z$ resuelven el mismo problema de Cauchy y por tanto
coinciden allí (unicidad incondicional de la pregunta 9), de modo
que un entorno de $t_1$ está en $Z$. Un subconjunto no vacío del
intervalo $J$ que sea a la vez \omterm{def:b2:metric:topology}{abierto} y cerrado en $J$ es todo
$J$ (los intervalos son \omterm{def:b2:metric:connected}{conexos},
\cref{thm:b2:metric:connectedness}): $y = z$ sobre $J$.

\textbf{24.} Aquí $f(t, y) = \alpha(t)y + \beta(t)$ es
$L$-lipschitziana en $y$ sobre todo $\intcc AB \times \R$ con $L =
\sup\abs\alpha$ (finito, pues $\alpha$ es \omterm{def:b2:metric:continuity}{continua} sobre un
segmento), y no hace falta ninguna banda $\intcc{y_0 - b}{y_0 +
b}$: tómese $X = C(\intcc AB)$ entero, sobre el cual $T$ está bien
definida. La inducción de la pregunta 8 se repite literalmente y
da $d_\infty(T^ny, T^nz) \leq \frac{(L(B - A))^n}{n!}d_\infty(y,
z)$: alguna iterada es una contracción sea cual sea la longitud $B
- A$, y la pregunta 4 concluye: una y solo una solución sobre todo
el segmento. La linealidad interviene exactamente una vez: hace
global en $y$ la cota \omterm{def:b2:metric:continuity}{lipschitziana}, eliminando el rectángulo y su
condición $h \leq b/M$.

\textbf{25.} La \emph{\omterm{def:b2:metric:complete}{completitud}} convirtió la sucesión de Cauchy
de iteradas en una solución de verdad (preguntas 1, 6 y 8), y su
ausencia dejó escapar $\sqrt2$ de $\Q$ (pregunta 13). La
\emph{contracción} dio la unicidad, el algoritmo y las barras de
error (preguntas 7 y 14). El \emph{truco de la iterada} eliminó la
condición de pequeñez $Lh < 1$, de modo que el intervalo
certificado solo depende de $M$ y no de $L$, y volvió globales las
ecuaciones lineales (preguntas 8 y 24). La \emph{\omterm{def:b2:metric:connected}{conexidad}}
ascendió la unicidad local a unicidad global (pregunta 23). Los
contraejemplos: $2\sqrt{\abs y}$ custodia la hipótesis
\omterm{def:b2:metric:continuity}{lipschitziana} (pregunta 11), $y^2$ custodia la localidad
(pregunta 12) y $\Q$ custodia la \omterm{def:b2:metric:complete}{completitud} (pregunta 13). La
cumbre es el teorema de Picard--Lindelöf (pregunta 8); cuando $f$
es meramente \omterm{def:b2:metric:continuity}{continua}, la existencia sobrevive pero la unicidad
muere, y la demostración cambia la contracción por la \omterm{def:b2:metric:compact}{compacidad}
de conjuntos de funciones: el teorema de Peano vía Arzelà--Ascoli,
en el volumen del tercer año.
\end{solution}
